% ==============================================================================
% ABSTRACT (ENGLISH)
% ==============================================================================
\chapter*{Abstract}
\addcontentsline{toc}{chapter}{Abstract}

\begin{tcolorbox}[
  enhanced,
  colback=vertprincipal!5!white,
  colframe=vertprincipal,
  boxrule=1.5pt, arc=6pt,
  left=10pt, right=10pt, top=8pt, bottom=8pt
]

This thesis presents the design and development of a complete automated
sentiment analysis system applied to student feedback on university courses.
Faced with the challenge of manually processing thousands of student reviews,
this project proposes an intelligent, automated, and real-time solution
leveraging recent advances in Natural Language Processing (NLP).

\vspace{0.5em}

The system is built upon the \textbf{CamemBERT} language model
(\textit{camembert-base}), a French variant of BERT developed by Inria and
CNRS, fine-tuned for three-class sentiment classification: \textbf{Negative},
\textbf{Neutral}, and \textbf{Positive}. The model was trained on a dataset
of 225 French student comments, balanced across the three classes, using the
following hyperparameters: learning rate of $2 \times 10^{-5}$, batch size
of 16, 3 training epochs, and maximum sequence length of 128 tokens. The
AdamW optimizer with a linear warmup scheduler was employed to ensure
stable convergence. The model achieves an accuracy of \textbf{79.41\%} on
the test set.

\vspace{0.5em}

The overall system architecture consists of four main components:
(1) a \textbf{RESTful backend} developed with FastAPI and SQLAlchemy
interfacing with a PostgreSQL database for result persistence;
(2) an \textbf{analytics dashboard} built with Streamlit and enriched
with interactive Plotly visualizations for real-time trend monitoring;
(3) a \textbf{mobile application} built with Flutter providing students
with a modern, intuitive interface to submit comments and view analysis
results; (4) a complete \textbf{NLP preprocessing pipeline} including
Unicode normalization, special character removal, and CamemBERT
tokenization.

\vspace{0.5em}

The obtained results demonstrate the feasibility and effectiveness of such
an approach, even with a limited dataset, paving the way for significant
improvements with additional annotated data.

\end{tcolorbox}

\vspace{0.8em}
\noindent\textbf{Keywords:} Sentiment analysis, NLP, CamemBERT, Machine
learning, FastAPI, Flutter, PostgreSQL, Text classification, Fine-tuning,
Deep learning, Transfer learning.

\cleardoublepage
